\lysh{14665}{4}

\begin{alist}
\item Οι τετμημένες των σημείων $A, B$ είναι οι λύσεις της εξίσωσης $f(x) = g(x)$ ισοδύναμα
$x^2 = x^3$, δηλαδή $x^2 - x^3 = 0$, οπότε $x^2(1 - x) = 0$ και άρα $x = 0 \text{ ή } x = 1$.
Για $x = 0$ είναι $f(0) = 0$ οπότε $A(0, 0)$.
Για $x = 1$ είναι $f(1) = 1$ οπότε $B(1, 1)$.

\item Από το σχήμα βλέπουμε ότι για κάθε $x \in (0, 1)$, δηλαδή για τις τετμημένες των
σημείων μεταξύ των $A$ και $B$, η γραφική παράσταση της $g$ είναι κάτω από τη
γραφική παράσταση της $f$, πράγμα που σημαίνει ότι $x^3 < x^2$.

Εναλλακτικά για κάθε $x \in (0, 1)$ είναι $x^2 > 0$ και $0 < x < 1$ οπότε $x \cdot x^2 < 1 \cdot x^2$ δηλαδή
$x^3 < x^2$.

\item Με βάση τα παραπάνω για κάθε $x \in (0, 1)$ είναι $x^3 < x^2$, π.χ. $\left(\dfrac{1}{2}\right)^3 < \left(\dfrac{1}{2}\right)^2$ αφού
ισοδύναμα $\dfrac{1}{8} < \dfrac{1}{4}$. Συνεπώς δεν είναι ο κύβος οποιουδήποτε αριθμού μεγαλύτερος
από το τετράγωνό του.

\item 
\begin{rlist}
\item Αφού $3 < \pi < 4$ είναι $3 - 3 < \pi - 3 < 4 - 3$ δηλαδή $0 < \pi - 3 < 1$ οπότε με βάση το 
$\beta)$ ερώτημα είναι $(\pi - 3)^3 < (\pi - 3)^2$.
\item Είναι
\[(\pi - 3)^3 < (\pi - 3)^2 \Leftrightarrow \pi^3 - 9\pi^2 + 27\pi - 27 < \pi^2 - 6\pi + 9 \Leftrightarrow \pi^3 - 10\pi^2 + 33\pi - 36 < 0\]
\end{rlist}
\end{alist}

